\id{ҒТАМР 65.09.03}{}

\begin{header}
\swa{М.Ж~Султанова, У.З~Сагындыков, А.С.~Сәдуақас, М.А.~Якияева, А.С~Камали, Н.~Акжанов}{СОКСЛЕТ ӘДІСІ АРҚЫЛЫ ЮГЛОН АЛУ ТЕХНОЛОГИЯСЫН ЖАСАУ ЖӘНЕ ӨНІМНІҢ ТҰРАҚТЫЛЫҒЫН БАҒАЛАУ}

\tsp{1}М.Ж~Султанова,
\tsp{3}У.З~Сагындыков,
\tsp{2}А.С.~Сәдуақас,
\tsp{1}М.А.~Якияева,
\tsp{2}А.С~Камали,
\tsp{3}Н.~Акжанов\envelope
\end{header}

\begin{affil}
\tsp{1}Алматы технологиялық университеті АҚ, Алматы, Қазахстан,

\tsp{2}Қазақ қайта өңдеу және тағам өнеркәсіптері ҒЗИ» ЖШС, Астана, Қазақстан,

\tsp{3}Л. Н. Гумилев атындағы Еуразия ұлттық университеті, Астана, Қазахстан

\corrauthor{Корреспондент-автор: e-mail: nurtore0308@gmail.com}
\end{affil}

Мақалада грек жаңғағының (Juglans regia L.) жасыл қабығы мен қабығы
биологиялық белсенді заттардың мол көзі ретінде зерттелді. Зерттеудің
негізгі мақсаты -- юглон (5-гидрокси-1,4-нафто\-хинон), яғни табиғи
химиялық қосылыстың концентрациясын бағалау, сондай-ақ олардың
экстракциялық шығымын арттыруға бағытталған тиімді технологияны әзірлеу
болып табылды. Жалпы шикізат Алматы облысында жемістердің нағыз
физиологиялық піскен кезеңінде жиналып, ≤\,45\,°C температурада
кептіріліп, герметикалық қара контейнерлерде сақталды. Дайындау кезеңі
екі сатылы ұнтақтауды қамтыды, гранулометриялық стандарттау арқылы
«қатты фазаның -- еріткішпен» беткі ауданы жоғары деңгейге жеткізілді,
бұл экстракция тиімділігін арттыруға мүмкіндік берді.

Биологиялық белсенді қосылыстарды экстракциялау Сокслет әдісімен,
«АСВ-6» жартылай автоматтандырылған аппаратында жүргізілді. Процесс
40--80\,°C аралығында орындалып, экстракт вакуумдық жағдайда ≤\,40\,°C
температурада концертрация жүргізілді. Юглон мөлшері жоғары тиімділікпен
сұйықтық хроматографиясы (ЖТШХ) әдісімен анықталды, ал әртүрлі шикізат
фракцияларынан алынған шығымның нақты мәндері калибреу қисығы негізінде
есептелді.

Зерттеу нәтижелер көрсеткендей, жасыл қабық юглонның ең бай көзі болып,
қабыққа қарағанда оның антиоксиданттық белсенділігі айтарлықтай жоғары
болып саналды. Юглонның шығымы бөлшектердің өлшеміне, экстракция
температурасына және уақытқа байланысты ауысты. Алынған тәжірибелік
деректер жаңғақ қалдықтарын биологиялық белсенді заттар көзі ретінде
тиімді пайдалануға, сондай-ақ фармацевтика, косметика салаларында,
функционалды тағам және қайта өңдеу өндірісінде қолдануға нақты ғылыми
негіз бола алады.

Зерттеу нәтижелері юглонның экстракциялық технологиясын жетілдіруге және
алынған экстракттің тұрақтылығы мен биологиялық белсенділігін сақтап
тұруға мүмкіндік береді. Алынған деректер жаңа биобелсенді өнімдер
әзірлеу және өндірістік ауқымда қолдану үшін маңызды практикалық және
ғылыми мәнге ие болып келеді.

{\bfseries Түйін сөздер:} грек жаңғағы, юглон, экстракт, технология,
Juglans regia L., биологиялық белсенді қосылыс, жасыл қабық.

\begin{header}
РАЗРАБОТКА ТЕХНОЛОГИИ ПОЛУЧЕНИЯ ЮГЛОНА МЕТОДОМ СОКСЛЕТА И ОЦЕНКА УСТОЙЧИВОСТИ ПРОДУКЦИИ

\tsp{1}М.Ж.~Султанова,
\tsp{3}У.З.~Сагындыков,
\tsp{2}А.С.~Сәдуақас,
\tsp{1}М.А Якияева,
\tsp{2}А.С.~Камали,
\tsp{3}Н.~Акжанов\envelope
\end{header}

\begin{affil}
\tsp{1}АО Алматинский технологический университет, Алматы, Казахстан,

\tsp{2}Казахский научно - исследовательский институт перерабатывающей и пищевой промышленности, Астана, Казахстан,

\tsp{3}Евразийский национальный университет им. Л.Н.Гумилева, Астана, Казахстан,

e-mail: nurtore0308@gmail.com
\end{affil}

В статье грецкий орех (Juglans regia L.) зеленая кожура и скорлупа были
изучены как обильные источники биологически активных веществ. Основной
целью исследования являлась оценка концентрации юглона
(5-гидрокси-1,4-нафтохинона), то есть природного химического соединения,
а также разработка эффективной технологии, направленной на повышение их
экстракционного выхода. Общее сырье собирали в Алматинской области в
период истинной физиологической спелости плодов, сушили при температуре
≤ 45 °C и хранили в герметичных черных контейнерах. Этап подготовки
включал двухступенчатое измельчение, при гранулометрической
стандартизации площадь поверхности «твердой фазы -- растворителем» была
доведена до высокого уровня, что позволило повысить эффективность
экстракции.

Экстракция биологически активных соединений производилась методом
Сокслета на полуавтоматическом аппарате "АСВ-6". Процесс выполнялся при
температуре от 40 до 80 °C, а экстракция проводилась при температуре ≤
40 °C в вакуумных условиях. Содержание юглона определяли методом
жидкостной хроматографии (ВЭЖХ) с высокой эффективностью, а фактические
значения расхода от различных сырьевых фракций рассчитывали на основе
кривой калибровки.

Исследование показало, что зеленая оболочка была самым богатым
источником юглона и считалась значительно более антиоксидантной
активностью, чем кора. Выход юглона менялся в зависимости от размера
частиц, температуры экстракции и времени. Полученные экспериментальные
данные могут служить реальной научной основой для эффективного
использования отходов грецкого ореха в качестве источника биологически
активных веществ, а также для использования в фармацевтической,
косметической промышленности, в производстве функциональных продуктов
питания и переработки.

Результаты исследования позволят усовершенствовать технологию экстракции
юглона и сохранить стабильность и биологическую активность полученного
экстракта. Полученные данные приобретают важное практическое и научное
значение для разработки и применения новых биоактивных продуктов в
промышленных масштабах.

{\bfseries Ключевые слова:} грецкий орех, юглон, экстракт, технология,
Juglans regia L., биологически активное соединение, зеленая оболочка.

\begin{header}
DEVELOPMENT OF TECHNOLOGY FOR THE PRODUCTION OF JUGLON BY THE SOXLET METHOD AND ASSESSMENT OF PRODUCT SUSTAINABILITY

\tsp{1}M.Zh.~Sultanova,
\tsp{3}U.Z.~Sagyndykov,
\tsp{2}A.S.~Saduakas,
\tsp{1}M.~AYakiyaeva,
\tsp{2}A.S.~Kamali,
\tsp{3}N.~Akzhanov\envelope
\end{header}

\begin{affil}
\tsp{1}Almaty Technological University JSC, Almaty, Kazakhstan,

\tsp{2}Kazakh Scientific Research Institute of Processing and Food Industry, Astana, Kazakhstan,

\tsp{3}L.N. Gumilyov Eurasian National University, Astana, Kazakhstan,

e-mail: nurtore0308@gmail.com
\end{affil}

Walnut (Juglans regia L.) green peel and shell have been studied as
abundant sources of biologically active substances. The main purpose of
the study was to evaluate the concentration of juglon
(5-hydroxy-1,4-naphthoquinone), a natural chemical compound, as well as
to develop an effective technology aimed at increasing their extraction
yield. The total raw materials were collected in the Almaty region
during the period of true physiological ripeness of the fruits, dried at
a temperature of 45 ° C and stored in airtight black containers. The
preparation stage included two--stage grinding, with granulometric
standardization, the surface area of the "solid phase - solvent" was
brought to a high level, which increased the extraction efficiency.

The extraction of biologically active compounds was carried out by the
Soxlet method on a semi-automatic device "ASV-6". The process was
carried out at a temperature of 40 to 80 °C, and the extraction was
carried out at a temperature of 40 °C under vacuum conditions. The
yuglon content was determined by liquid chromatography (HPLC) with high
efficiency, and the actual flow rates from various raw material
fractions were calculated based on the calibration curve.

The study showed that the green shell was the richest source of juglon
and was considered to have significantly more antioxidant activity than
the bark. The yield of yuglon varied depending on the particle size,
extraction temperature, and time. The experimental data obtained can
serve as a real scientific basis for the effective use of walnut waste
as a source of biologically active substances, as well as for use in the
pharmaceutical, cosmetic, and functional food production and processing
industries.

The results of the study will make it possible to improve the technology
of yuglon extraction and maintain the stability and biological activity
of the resulting extract. The data obtained acquire important practical
and scientific significance for the development and application of new
bioactive products on an industrial scale.

{\bfseries Keywords}: walnut, juglon, extract, technology, Juglans regia
L., biologically active compound, green shell.

\begin{multicols}{2}
{\bfseries Кіріспе}. Juglans mandshurica ағаштарының қалдықтары биологиялық
белсенді қосылыстардың, атап айтқанда юглонның әлеуетті көзі болып
табылады. Li C. және оның әріптестері осы қалдық бұтақтардан юглонды
бөліп алу мақсатында «су майда» (MBMAE) микроэмульсиясына негізделген
микротолқынды экстракция әдісін ұсынды. Массалық қатынасы 27:13,5:4,5:55
болатын микроэмульсия құрамын (Твин 80, n-пропанол, n-гексан және су) pH
5,6 және 40\,°C температурада жетілдіру арқылы қосылысты 63 секунд
ішінде тиімді алу мүмкіндігі қамтамасыз етілді. Юглонның шығымы
4,58\,мг/г болды, бұл дәстүрлі экстракция әдістеріне қарағанда
1,86--6,65 есе жоғары тиімділікті көрсетті. Осыған ұқсас тиімділіктің
артуы Ramezani N. және әріптестері жүргізген зерттеуде де байқалды, онда
суперкритикалық CO₂-экстракция этанолды қос еріткіш ретінде колданылып
жүзеге асырылды. Зерттеу нәтижелері көрсеткендей, қысым бөлінген юглон
мөлшерін анықтайтын негізгі фактор болып шықты, ал температураның
көтерілуі және бөлшек өлшемінің ұлғаюы шығымды төмендеткенін жазған. Ең
көп юглон шығымы 37,26\,мг деңгейінде тіркелді, бұл 150\,бар қысым,
35\,°C температура және бөлшек өлшемі 375\,мкм болған жағдайда алынды.
Бұл деректер суперкритикалық экстракцияның зертханалық жағдайда және
өндірістік қолданудағы практикалық маңыздылығын көрсетеді {[}1,2{]}.

Сонымен қатар зерттеушілер грек жаңғағының әртүрлі бөліктеріндегі
фенолды қосылыстар мен антиоксиданттық белсенділікке назар аударды.
Jahanban-Esfahlan A. және әріптестері көрсеткендей, қабық, тұқым қабығы
және жапырақтар фенолдармен қанық және жоғары антиоксиданттық
автивтілікке ие, тіпті дәстүрлі түрде қоректік құндылық тек жаңғақ
ядросымен байланысты болғанымен, Luksiene J. және әріптестері
жемістердің орталық перделері фенолды қосылыстар мен
проантоцианидиндердің зор көзі екенін көрсетті: фенолдардың жалпы
мөлшері құрғақ массаның 1 грамында 131,55-тен 530,92\,мг галлов қышқылы
эквивалентінде, ал проантоцианидиндер - 1,14-тен 7,65\,мг (-)-эпикатехин
эквивалентінде өзгерді. Биоактивті қосылыстардың ең жоғары
концентрациясы Шяуляй үлгісінде тіркелді, бұл олардың табиғи
нутрицевтикалық және антиоксиданттық өнімдер өндірісінде жоғары
мүмкіндігін растайды {[}3,4{]}.

Ғалымдар жапырақтар мен Juglans гибридтеріне, мысалы, Juglans major 209
× Juglans regia, фенолдардың айқын антиоксиданттық белсенділікке ие көзі
ретінде қызығушылық танытты. Fernández-Agulló A. және әріптестері
көрсеткендей, экстракцияның оңтайды жағдайлары (75\,°C, 120\,мин, 50\,\%
этанол) 30,17\,\% шығыммен экстракт алуға мүмкіндік беріп, құрамында
кверцетин 3-β-D-глюкозид, неохлороген қышқылы және хлороген қышқылы
анықталды. Сонымен қоса, Pop A. және әріптестері Juglans regia L. ерлер
гүлдерін ультрадыбыстық экстракция арқылы зерттеп, ацетонның
полифенолдар мен токоферолдарды алу үшін ең тиімді еріткіш екенін
анықтады. Алынған экстрактар жоғары антиоксиданттық белсенділікті
көрсетті, α-глюкозидаза және тирозиназа ферменттік белсенділігін тежеді,
сонымен қатар жасушаларға зиян келтірмей қатерлі ісікке қарсы төтеп
берді {[}5,6{]}.

Грек жаңғағының қабығы агроөнеркәсіптік қалдық ретінде кең таралған,
сондай-ақ биологиялық белсенді қосылыстардың көзі болып табылады. Seabra
I. J. және әріптестері жоғары қысым кезінде екі сатылы экстракция
қолданды: алдымен суперкритикалық CO₂, содан кейін субкритикалық
CO₂/этанол/су қоспасы температурасы 35--45\,°C және қысымы 10--30\,МПа
болатын жағдайда зерттеген. ГХ-МС талдау алкандарды, терпеноидтарды және
алкалоидтарды анықтады, ал юглон мен 1,4-нафтохинонның мөлшері 0,33-тен
2,77\,мг/г диапазонында болды. Екінші сатыдағы экстракциядағы
фенолдардың жалпы деңгейі 71,5--116,8\,мг GAE/г құрады. Юглон мөлшеріне
қарамастан экстрактардың жоғары антиоксиданттық белсенділігі жоғары
қысым тәсілдерінің фармацевтика, ауыл шаруашылығы және тамақ өнеркәсібі
үшін селективті және биологиялық құнды экстракттар алу мүмкіндігін
растайды {[}7{]}.

Грек жаңғағының (Juglans spp.) әртүрлі бөліктерінен биологиялық белсенді
қосылыстарды экстракциялау зерттеушілердің қызығушылығын оятты, өйткені
оларда фенолды қосылыстар, флавоноидтар және юглон көп мөлшерде
кездеседі, және олар антиоксиданттық, микробқа қарсы және цитотоксикалық
қасиеттерге ие. Мысалы, Rajković K. M. зерттеуінде Juglans nigra
жапырақтарынан 50\,\% этанол және еріткіш/шикізат қатынасы 20\,кг/кг
қолдану арқылы экстракцияның оңтайлы жағдайлары фенолды қосылыстардың,
соның ішінде 3-O-кофеоилхин қышқылы, кверцетин-3-O-галактозид және
кверцетин-3-O-рамнозидтің жоғары шығымын қамтамасыз еткені, сондай-ақ
макро- және микроэлементтердің сақталғаны көрсетілді {[}8{]}. Zhang Y.
G. зерттеулерінде жеміс перделері, қабық және гүлдердегі фенолды
қосылыстар анықталып, 50 қосылысты тіркеді; қабықта ең көп таниндер
болғаны анықталса, ал антиоксиданттық белсенділік әсіресе жоғары болды,
бұл осы бөліктерді функционалды өнімдер өндіруде қолданудың
перспективалы екенін растайды {[}9{]}. Rusu M. E. көрсеткендей, жеміс
перделері полифенолды қосылыстар мен конденсацияланған таниндерге бай,
ал тирозиназа ингибиторлық тесттері олардың тағамдық, фармацевтикалық
және косметикалық өнеркәсіп үшін құндылығын дәлелдеді {[}10{]}.

Juglans regia жапырақтары да құнды биологиялық белсенді қосылыстардың
көзі болып табылады. Cordova I. W. суда ерітетін эвтектикалық
еріткіштерді пайдалану арқылы экстракцияның экологиялық таза әдісін
жасап, жоғары фенол шығымын қамтамасыз етіп, айқын антиоксиданттық және
микробқа қарсы белсенділікті көрсетті {[}11{]}. Paşayeva L.
көрсеткендей, жапырақтар мен жасыл қабық антихолинэстераза
белсенділігіне, антиоксиданттық және цитотоксикалық қасиеттерге ие, ал
негізгі белсенді компоненттер ретінде эллаг қышқылы мен хлороген қышқылы
табылды {[}12{]}. Julia B. зерттеуінде жеміс пелликулдарын
ультрадыбыстық өңдеу полифенолды қосылыстар мен дубильді заттардың
шығымын арттырған, алайда өңдеу параметрлері әрбір сұрыпқа мұқият
таңдалуы тиіс; мысалы, Valentina' s Gift сұрыбы үшін
оптималды жағдайлар биологиялық белсенді заттардың көп мөлшердегі
шығымын қамтамасыз етті {[}13{]}. Осы деректердің жиынтығы Juglans
ағаштарының әртүрлі бөліктерінен айқын биологиялық белсенділігі бар
экстракттер алуға және оларды тағам, қайта өңдеу, косметика және
фармацевтика өнеркәсібінде қолдануға жоғары мүмкіндік бар екенін
дәлелдейді.

{\bfseries Материалдар мен әдістер.} Зерттеу объектісі ретінде Алматы
облысында жемістердің нағыз физиологиялық піскен кезеңінде (2025 жылғы
шілде--тамыз айларында) жиналған Juglans regia L. жаңғағының жасыл
қабығы мен қабығы алынды. Шикізатты фотодеструкцияны және юглон мен
басқа биологиялық белсенді қосылыстардың ыдырауын болдырмау мақсатында
тікелей күн сәулесінен қорғай отырып, ≤\,45\,°C температурада тұрақты
массаға дейін кептірді. Кептірілген материал химиялық тұрақтылықты
сақтау мақсатында герметикалық қара контейнерлерде ≤\,25\,°C
температурада және ≤\,12\,\% салыстырмалы ылғалдылықта сақталды.

Шикізатты дайындау екі сатылы ұнтақтауды қамтамасыз етті: алдымен
Novital Magnum 4V соққылы ұнтақтағышпен 1--2,5\,мм фракциясына дейін
алдын ала ұнтақталды, кейін МШЛ-1П лабораториялық шарлы ұнтақтағышпен
0,15\,мм бөлшектерге дейін майдаланды. Ұнтақтың гранулометриялық құрамын
стандарттау үшін олар қосымша 0,1--0,25\,мм лабораториялық електерден
өткізілді. Мұндай дайындық «қатты фазаның -- еріткішпен» байланысатын
беткі ауданын барынша арттыруға мүмкіндік беріп, экстракцияның
тиімділігін айтарлықтай жоғарылатты.

Биологиялық белсенді қосылыстарды экстракциялау Сокслет әдісімен,
«АСВ-6» жартылай автоматтандырылған аппаратындағы колбаны қолдану арқылы
жүргізілді. Алдын ала ұнтақталған шикізаттың (2,00\,±\,0,01\,г)
ілмектері экстракциялық гильзаға орналастырылды, ал колбаға 45\,мл
еріткіш (этанол немесе этилацетат) құйылды. Үрдіс 40--80\,°C
температурада жүргізілді: еріткішпен байланыс кезеңі -- 30\,минут, таза
еріткішпен шаю -- 60--180\,минут. Экстракциядан кейін алынған экстракт
салқындатылып, сүзгіден өткізілді және вакуумдық жағдайда ≤\,40\,°C
температурада 10\,мл көлеміне дейін буланды. Алынған концентраттар юглон
мен басқа фенолды қосылыстардың тұрақтылығын сақтау мақсатында қараңғы
флакондарда t ≤\,8\,°C температурада сақталды.

Тәжірибеде алынған экстрактардағы юглон мөлшері жоғары тиімділікпен
сұйықтық хроматографиясы (ЖТШХ) әдісімен анықталды. Талдау үшін
стандартталған юглон ерітінділері, C18 колонкасы және 254\,нм толқын
ұзындығындағы УФ-спектроскопиялық детектор ұсынылды. Сандық көрсеткіштер
калибрлеу қисығына сәйкес есептелді, бұл әртүрлі шикізат фракцияларынан
алынған юглон шығымының дәл мәндерін алуға мүмкіндік берді.

{\bfseries Нәтижелер мен талқылау.} Грек жаңғағының (Juglans regia L.)
жасыл қабығы мен қабығының биологиялық белсенді қосылыстар көзі ретінде
потенциалын бағалау үшін олардың химиялық құрамына кешенді талдау
жүргізілді. Арнайы назар нафтохинон қосылыстарына, ең алдымен юглонға
(5-гидрокси-1,4-нафтохинон) назар аударылды, себебі ол биологиялық
белсенді маркер ретінде қызмет етеді.

Зерттеу нәтижелері көрсеткендей, юглон мен басқа фенолды қосылыстардың
концентрациясы жаңғақтың морфологиялық бөліктеріне қарай айтарлықтай
өзгереді. Жасыл қабық фенолдардың мөлшері бойынша қабыққа қарағанда
әлдеқайда бай болып келсе, бұл метаболиттердің жемістің сыртқы беткі
тіндерінде орналасуы және патогендерге сезімталдығына байланысты
атқаратын қорғаныштық функциясымен түсіндіріледі. (1-кесте).
\end{multicols}

\tcap{1 - кесте. Juglans regia L. өсімдік шикізатының химиялық сипаттамасы}
% \usepackage{tabularray}
\begin{longtblr}[
  label = none,
  entry = none,
]{
  cells = {c},
  cells = {font = \small},
  row{1} = {font=\bfseries\small},
  hlines,
  vlines,
}
Көрсеткіш & Жасыл қабық & Қабық\\
Ылғалдылық, \% & 8,5 ± 0,4 & 6,2 ± 0,3\\
Күл, \% & 7,1 ± 0,2 & 4,8 ± 0,2\\
Жалпы фенолды қосылыстар, мг ГҚ/г & 85–95 & 28–35\\
Юглон, мг/г құрғақ зат & 5,8–6,5 & 1,2–1,6
\end{longtblr}

\begin{multicols}{2}
Кесте деректерінен сүйенсек, грек жаңғағының жасыл қабығында фенолды
қосылыстардың (85--95 мг ГҚ/г) және юглонның (5,8--6,5 мг/г) мөлшері
қабықшаға қарағанда едәуір жоғары. Бұл құбылыс аталған қосылыстардың
жемістің сыртқы тіндерінде атқаратын қорғаныштық қасиетімен байланысты.
Жасыл қабықтың ылғалдылығы мен күлділігі төмен деңгейде сақталады, бұл
шикізатты сақтау және технологиялық өңдеу кезінде оның тұрақтылығын
қамтамасыз етеді. Осыған байланысты грек жаңғағының жасыл қабығы
биологиялық белсенді заттарды экстракциялау және олардың антиоксиданттық
әлеуетін бағалау үшін ең қолайлы шикізат ретінде пайдаланылады.

Экстракция процесінің тиімділігін қамтамасыз ету үшін шикізатқа алдын
ала дайындық кезеңі жүргізіледі. Бұл кезеңде қабықша мен жасыл қабық
жақсылап іріктеліп, қажетті өлшемге дейін ұнтақталады және оптималды
ылғалдылық деңгейіне кептіріледі. Дайындалған шикізаттың біркелкілігі
және құрылымдық сипаттамалары юглонның ерітіндіге максималды өтуіне
ықпал етеді, яғни массалық тасымалдау мен диффузия процестерін
оңтайландырады (сур.1).
\end{multicols}

\fig[0.5\textwidth]{p2/image2}[1-сурет. Грек жаңғағының (Juglans regia L.) өсімдік шикізатын
юглон экстракциясына дайындау: қабықша және жасыл қабық»]
\fig[0.6\textwidth]{p2/image19}[2-сурет. Экстракция кезінде бөлшек мөлшерінің юглон шығымына әсері]

\begin{multicols}{2}
Суретте көрсетілгендей, жасыл қабық пен қабықшаны дайындау барысында
бөлшектердің өлшемі, біркелкілігі және материалдың басты критерий болып
табылады. Бұл дайындық юглонның кейінгі экстракциядағы максималды
шығымын қамтамасыз етуге және алынған экстрактінің сапасын
тұрақтандыруға негіз болады.

Алынған эксперименттік деректер шикізаттың гранулометриялық құрамының
юглонды экстракциялау үрдісіндегі айқындаушы рөлін көрсетеді. Диаграмма
деректері бойынша бөлшек өлшемінің 0,05-тен 0,20 мм-ге дейін көбеюі
юглон шығымының 3,9-дан 6,0 мг/г-ге дейін артуымен қатар жүреді (сур.2).

Бөлшек өлшемінің одан әрі 0,25--0,30 мм-ге дейін көбеюі юглон шығымының
5,2--4,7 мг/г деңгейіне төмендеуіне әкеледі. Мұндай үрдіс экстрагирлеуші
ортамен жанасу кезінде бетінің азаюына және биологиялық белсенді
қосылыстың диффузиялық тасымалының баяулауына байланысты деп пайымдауға
болады. Осылайша, шамамен 0,20 мм бөлшек өлшемі юглонды кейінгі
экстракциялау үшін ең ұтымды гранулометриялық көрсеткіш ретінде
қарастырылуы мүмкін.

Тәжірибелік нәтижелерді талдау экстракция температурасының жоғарылауы
юглон шығымының біртіндеп артуымен сипатталатынын көрсетті. Атап
айтқанда, 40 °C температурада юглонның шығымы 3,65 мг/г болса, 60--80 °C
аралығында бұл көрсеткіш 4,85--5,05 мг/г деңгейіне дейін көбейді
(2-кесте).
\end{multicols}

\tcap{2 - кесте. Экстракция үрдісінің температуралық аспектілері}
\begin{longtblr}[
  label = none,
  entry = none,
]{
  cells = {c},
  cells = {font = \small},
  row{1} = {font=\bfseries\small},
  hlines,
  vlines,
}
Экстракция температурасы, °C & Юглон шығымы, мг/г\\
40 & 3,65\\
60 & 4,85\\
70 & 4,95\\
80 & 5,05
\end{longtblr}

\fig[0.6\textwidth]{p2/image20}[1 - сурет. Экстракция уақыты және экстракттағы юглонның жинақталу динамикасы]

\begin{multicols}{2}
Температураның жоғарылауына байланысты юглон шығымының артуы масса
алмасу үрдістерінің қарқындануымен, қосылыстың еріткіште жақсырақ
еруімен және өсімдік матрицасынан диффузия жылдамдығының жоғарылауымен
түсіндіріледі. Сонымен қатар, зерттелген температуралық аралықта
юглонның термиялық ыдырау белгілері байқалмады, бұл экстракцияның
берілген шарттарында қосылыстың жеткілікті термотұрақтылығын көрсетеді.

Юглонның экстрактта жинақталу динамикасы экстракция уақытының ұлғаюына
қарай оның концентрациясының тұрақты түрде артуымен сипатталады.
Экстракцияның алғашқы 30--90 минуты аралығында юглонның қарқынды
бөлініуі байқалып, шығым мөлшері 3,65-тен 4,75 мг/г-ге дейін артты
(сур.3).

Экстракция ұзақтығы 120 және 180 минутқа дейін ұзартылған кезде юглонның
шығымы сәйкесінше 5,05 және 6,05\,мг/г мәндеріне дейін жетті. Зерттелген
уақыт аралығында айқын өсінің байқалмауы юглонды бөліп алу процесінің
негізінен диффузиялық механизммен шектелетінін көрсетеді. Яғни,
экстракция уақыты ұлғайған сайын, юглон өсімдік шикізатынан ерітіндіге
біртіндеп және толық өткен кезде, бұл алынған экстракттің
концентрациясының тұрақтылығын қамтамасыз етеді.

Алынған экстракттің қасиеттерін бағалау үшін оның негізгі
физико-химиялық көрсеткіштері қарастырылды. Юглонның молекулалық
құрылымы мен физико-химиялық сипаттамалары экстракция үрдісінің
тиімділігін түсінуге мүмкіндік береді, сонымен қатар алынған өнімнің
тұрақтылығы мен келесі қолданылуына баға береді (3-кесте).
\end{multicols}

\tcap{1 - кесте. Юглон экстрактісінің физика-химиялық қасиеттерінің көрсеткіші}
\begin{longtblr}[
  label = none,
  entry = none,
]{
  width = 0.8\linewidth,
  colspec = {Q[500]Q[500]},
  cells = {c},
  row{1} = {font=\bfseries},
  hlines,
  vlines,
}
Қасиеті & Мәні\\
Молекулалық массасы & 174,15 г/моль\\
Балқу температурасы & 161–163 °C\\
Суда ерігіштігі & Нашар ериді\\
Органикалық еріткіштерде ерігіштігі & Жақсы ериді (диоксан, ацетон, хлороформ, бензол, сірке қышқылы)\\
pKa & 6,96\\
Түсі & Сары-қызғылт кристалды ұнтақ\\
Ыдырау температурасы & 385 °C
\end{longtblr}

\begin{multicols}{2}
Юглонның молекулалық құрылымы, ерігіштік қасиеттері және температураға
төзімділігі оның экстракция тиімділігін, сондай-ақ концентрацияны
бақылау мен сақтау шарттарын анықтауға мүмкіндік береді. Мысалы, оның
суда нашар еруі мен органикалық еріткіштерде жақсы еруі этанол немесе
этилацетат сияқты полярлық қасиеті орташа еріткіштердің қолданылуын
айқындайды. Сонымен қатар, жоғары ыдырау температурасы (385\,°C)
юглонның термотұрақытлығын көрсетеді, бұл вакуумдық буландыру немесе
төмен температуралы сақтау кезінде биоактивтіліктің сақталуын қамтамасыз
етеді.

Юглонды Juglans regia L. қабығынан алу процесі оның молекулалық
қасиеттеріне және биобелсенділігін сақтауға бағытталған операциялық
блоктар арқылы ұйымдастырылған. Шикізатты дайындау, кептіру және
ұнтақтау кезеңдері экстракция тиімділігін қамтамасыз етеді, ал соңғы
концентрациялау мен сақтау процестері алынған қосылыстардың химиялық
тұрақтылығын сақтап тұруға мүмкіндік береді (сур.4).

Суретте көрсетілген технологиялық схема юглонды экстракциялау үрдісінің
әрбір кезеңін нақты көрсетеді. Шикізатты қабылдау және бастапқы
дайындау: механикалық қосылыстарды жойып, морфологиялық сұрыптауды
қамтамасыз етеді. Кептіру кезеңі: фотодеструкция мен термиялық ыдырауды
болдырмай, биологиялық белсенділікті сақтайды. Ұнтақтау және фракциялық
стандарттау: экстракция тиімділігін арттырады, ал Сокслет әдісімен
экстракциялау: юглон мен басқа фенолды қосылыстарды бөліп алудың жоғары
нәтижесін береді. Сүзу және вакуумдық концентрация: алынған экстракттан
еріткішті тиімді түрде бөліп, концентраттың сапасын қамтамасыз етеді.
Соңғы кезеңде: дайын концентрат қара ыдыстарда төмен температурада
сақталып, химиялық тұрақтылық пен биоактивтілікті сақтауға мүмкіндік
береді. Бұл сызба юглонның молекулалық қасиеттерін ескере отырып,
зертханалық жағдайда және өндірістік масштабта қолдануға ғылыми тұрғыдан
негізделген.
\end{multicols}

\fig{p2/image21}[1 - сурет. Юглон экстракциясын алу технологиясы]

\begin{multicols}{2}
{\bfseries Қорытынды.} Жүргізілген зерттеулер нәтижелері бойынша юглонды
экстракциялау үрдісінің кешенді сипаттамасын беріп, оның ғылыми негізін
нақтылады. Талдау көрсеткендей, морфологиялық бөліктер мен олардың
физика-химиялық қасиеттері, гранулометриялық құрылымы және ылғалдылық
экстракция тиімділігін анықтайтын негізгі көрсеткіштер болып табылады.
Ұнтақтау мен фракциялық стандарттау арқылы қатты фазаның беткі ауданы
артып, масса тасымалдау мен диффузия үрдістері жетілдіріледі, бұл юглон
мен басқа фенолды қосылыстардың ерітіндіге максималды өтуіне мүмкіндік
береді. Экстракция температурасы мен ұзақтығының таңдалуы қосылыстардың
шығымын арттырады, ал Сокслет әдісі биобелсенді компоненттердің толық
бөлінуін қамтамасыз етеді.

Сүзу және вакуумдық концентрация кезеңдері алынған экстракттың химиялық
тұрақтылығын сақтап, юглонның термотөзімділігін қамтамасыз етеді.
Нәтижесінде алынған өнім құрғақ, юглонның жоғары концентрациялы
экстракты мол көзі болып шықты, құрамындағы фенолды қосылыстар мен
биологиялық белсенді компоненттердің антиоксиданттық қасиеті толық
сақталған. Бұл экстракт фармацевтикалық, косметикалық және функционалды
тағамдық және қайта өңдеу өндірістерінде қолдануға дайын, әрі ғылыми
тұрғыдан негізделген, өндірістік қолдануға жарамды компонент болып
табылады. Жүргізілген жұмыс экстракция үрдісінің әрбір кезеңінің рөлін,
олардың өзара байланысын және алынған өнім сапасына әсерін терең
түсінуге мүмкіндік берді, сондай-ақ технологиялық үрдісті жетілдіруге
және болашақта өндірістік ауқымда шығаруға толық негізін қалады.

\emph{Қаржыландыру.} Жұмыс Қазақстан Республикасының Ғылым және жоғары
білім министрлігі, AP26197887 "Тамақ өнеркәсібінде қолдану үшін юглонды
өсімдік көздерінен экстракциялау технологиясын әзірлеу" ҒЗЖ
қаржыландыратын жоба шеңберінде орындалды.

Авторлар «Қазақ қайта өңдеу және тағам өнеркәсіптері ғылыми-зерттеу
институты» ЖШС (Астана, Қазақстан) қызметкерлеріне тәжірибелік
зерттеулерді жүргізуге және алынған нәтижелерді талқылауға көрсеткен
жәрдемі үшін шын жүректен алғыс білдіреді.
\end{multicols}

\begin{center}
{\bfseries Әдебиеттер}
\end{center}

\begin{refs}
1. Li C. et al. A novel method to extract juglone from Juglans
mandshurica waste branches using a water-in-oil microemulsion //Waste
and Biomass Valorization. -2022.-Vol.13(3). - P.1547-1563. DOI
10.1007/s12649-021-01611-x.

2. Ramezani N. et al. Juglone extraction from walnut (Juglans regia L.)
green husk by supercritical CO2: Process optimization using Taguchi
method //Journal of Environmental Chemical Engineering. - 2020.-
Vol.8(3). - P.103776. DOI10.1016/j.jece.2020.103776.

3. Jahanban-Esfahlan A. et al. A comparative review on the extraction,
antioxidant content and antioxidant potential of different parts of
walnut (Juglans regia L.) fruit and tree //Molecules. --
2019. -Vol.24(11). - P.2133. DOI 10.3390/molecules24112133.

4. Luksiene J. Et al. Examining the Role of Extraction Techniques and
Regional Variability in the Antioxidant and Phytochemical Composition of
Juglans regia L. Septa //Plants. -2025.- Vol.14(16)-P.2524. DOI
10.3390/plants14162524.

5. Fernández-Agulló A. et al. Optimization of the extraction of
bioactive compounds from walnut (Juglans major 209 x Juglans regia)
leaves: Antioxidant capacity and phenolic profile //Antioxidants. -2019.
- Vol.9(1). -- P.18. DOI 10.3390/antiox9010018.

6. Pop A. et al. Enhanced recovery of phenolic and tocopherolic
compounds from walnut (Juglans regia L.) male flowers based on process
optimization of ultrasonic assisted-extraction: Phytochemical profile
and biological activities //Antioxidants. -2021.- Vol.10(4). - P.607.
DOI 10.3390/antiox10040607.

7. Seabra I. J. et al. Two-step high pressure solvent extraction of
walnut (Juglans regia L.) husks: scCO2+ CO2/ethanol/H2O //Journal of CO2
Utilization. - 2019.- Vol.34.- P.375-385. DOI
10.1016/j.jcou.2019.07.028.

8. Rajković K. M. et al. Optimization of extraction yield and chemical
characterization of optimal extract from Juglans nigra L. leaves
//Chemical Engineering Research and Design. -2020.- Vol.157.- P.25-33.
DOI 10.1016/j.cherd.2020.03.002.

9. Zhang Y. G. et al. Comparison of phenolic compounds extracted from
Diaphragma juglandis fructus, walnut pellicle, and flowers of Juglans
regia using methanol, ultrasonic wave, and enzyme assisted-extraction
//Food chemistry. - 2020.- Vol.321. -- P.126672.
DOI 10.1016/j.foodchem.2020.126672.

10. Rusu M. E. et al. Process optimization for improved phenolic
compounds recovery from walnut (Juglans regia L.) septum: Phytochemical
profile and biological activities //Molecules. -2018. --Vol.23(11):
2814. DOI 10.3390/molecules23112814.

11. Cordova I. W. et al. Extraction of phenolic compounds from Juglans
regia L. leaves using aqueous solutions of eutectic solvents
//Separation and Purification Technology. -2025.- Vol.354(7): 129214.
DOI 10.1016/j.seppur.2024.129214.

12. Paşayeva L. Et al. Optimizing health benefits of walnut (Juglans
regia L.) agricultural by-products: impact of maceration and Soxhlet
extraction methods on phytochemical composition, enzyme Inhibition,
antioxidant, antimicrobial, and cytotoxic activities //Food Bioscience.
-2025. Vol.64:105923. DOI 10.1016/j.fbio.2025.105923.

13. Bazarniva J. et al. Extraction of polyphenolic compounds from the
Juglans regia L. pellicles of using ultrasound //The chemistry of plant
materials-2023. -№.1.- P.273-278. DOI 10.4258/jcprm.2023011970.
\end{refs}

\begin{info}
\hspace{1em}\emph{{\bfseries Авторлар туралы мәліметтер}}

Султанова М.Ж. - докторант, Алматы технологиялық университеті АҚ,
Алматы, Қазақстан, e-mail:
sultanova.2012@mail.ru;

Сагындыков У.З. - биология ғылымдарының кандидаты, қауымдастырылған
профессор, Л. Н. Гумилев атындағы Еуразия ұлттық университеті, Астана,
Қазахстан, e-mail: outemourate@list.ru;

Сәдуақас А.С. - ғылыми қызметкер, «Қазақ қайта өңдеу және тағам
өнеркәсіптері ҒЗИ» ЖШС, Астана, Қазақстан, e-mail: aykon96@mail.ru;

Якияева М.А. - PhD, қауымдастырылған профессор, Алматы технологиялық
университеті АҚ, Алматы, Қазақстан, e-mail: yamadina88@mail.ru;

Камали А.С. - кіші ғылыми қызметкер, «Қазақ қайта өңдеу және тағам
өнеркәсіптері ҒЗИ» ЖШС, Астана, Қазақстан, e-mail: akerke.kamali@bk.ru;

Акжанов Н. - докторант, Л. Н. Гумилев атындағы Еуразия ұлттық
университеті, Астана, Қазахстан, e-mail: nurtore0308@gmail.com.

\hspace{1em}\emph{{\bfseries Information about the authors}}

Sultanova M.Zh. - Doctoral student, Almaty Technological University JSC,
Almaty, Kazakhstan, e-mail: sultanova.2012@mail.ru;

Sagyndykov U.Z. - Candidate of Biological Sciences, Associate Professor,
L.N. Gumilyov Eurasian National University, Astana, Kazakhstan, e-mail:
outemourate@list.ru;

Saduakas A.S. - Researcher, Kazakh Scientific Research Institute of
Processing and Food Industry LLP, Astana, Kazakhstan, e-mail:
aykon96@mail.ru;

Yakiyaeva M.A.- Ph.D., Associate Professor, Almaty Technological
University, Almaty, Kazakhstan, e-mail: yamadina88@mail.ru;

Kamali A.S. - Junior Researcher, Kazakh Scientific Research Institute of
Processing and Food Industry LLP, Astana, Kazakhstan, e-mail:
akerke.kamali@bk.ru;

Akzhanov N.- Doctoral student, L.N. Gumilyov Eurasian National
University, Astana, Kazakhstan, e-mail:
nurtore0308@gmail.com.
\end{info}
